\documentclass[11pt]{article}
% \documentclass[11pt]{ctexart}



\usepackage{xcolor}
% Change "review" to "final" to generate the final (sometimes called camera-ready) version.
% Change to "preprint" to generate a non-anonymous version with page numbers.
\usepackage[final]{acl}

% Standard package includes
\usepackage{times}
\usepackage{latexsym}
\usepackage{amsmath}
\usepackage{booktabs}
% For proper rendering and hyphenation of words containing Latin characters (including in bib files)
\usepackage[T1]{fontenc}
% For Vietnamese characters
% \usepackage[T5]{fontenc}
% See https://www.latex-project.org/help/documentation/encguide.pdf for other character sets

% This assumes your files are encoded as UTF8
\usepackage[utf8]{inputenc}

% This is not strictly necessary, and may be commented out,
% but it will improve the layout of the manuscript,
% and will typically save some space.
\usepackage{microtype}

% This is also not strictly necessary, and may be commented out.
% However, it will improve the aesthetics of text in
% the typewriter font.
\usepackage{inconsolata}

%Including images in your LaTeX document requires adding
%additional package(s)
\usepackage{graphicx}

% For inline-coded plots (Figure 1: rounds distribution).
\usepackage{pgfplots}
\pgfplotsset{compat=1.18}
\usepgfplotslibrary{groupplots}
\definecolor{stdblue}{HTML}{4C72B0}
\definecolor{optorange}{HTML}{DD8452}

% If the title and author information does not fit in the area allocated, uncomment the following
%
%\setlength\titlebox{<dim>}
%
% and set <dim> to something 5cm or larger.

\title{Evaluating and Improving LLM-Based Agents in Long-Horizon OS Task Execution}

% Author information can be set in various styles:
% For several authors from the same institution:
% \author{Author 1 \and ... \and Author n \\
%         Address line \\ ... \\ Address line}
% if the names do not fit well on one line use
%         Author 1 \\ {\bf Author 2} \\ ... \\ {\bf Author n} \\
% For authors from different institutions:
% \author{Author 1 \\ Address line \\  ... \\ Address line
%         \And  ... \And
%         Author n \\ Address line \\ ... \\ Address line}
% To start a separate ``row'' of authors use \AND, as in
% \author{Author 1 \\ Address line \\  ... \\ Address line
%         \AND
%         Author 2 \\ Address line \\ ... \\ Address line \And
%         Author 3 \\ Address line \\ ... \\ Address line}

\author{ Hunter Zhang \And Guangwen Xiong \And Wanyi Chen \And Leyan Chen \And Jingyu Huang \And Yanting Guo}

%\author{
%  \textbf{First Author\textsuperscript{1}},
%  \textbf{Second Author\textsuperscript{1,2}},
%  \textbf{Third T. Author\textsuperscript{1}},
%  \textbf{Fourth Author\textsuperscript{1}},
%\\
%  \textbf{Fifth Author\textsuperscript{1,2}},
%  \textbf{Sixth Author\textsuperscript{1}},
%  \textbf{Seventh Author\textsuperscript{1}},
%  \textbf{Eighth Author \textsuperscript{1,2,3,4}},
%\\
%  \textbf{Ninth Author\textsuperscript{1}},
%  \textbf{Tenth Author\textsuperscript{1}},
%  \textbf{Eleventh E. Author\textsuperscript{1,2,3,4,5}},
%  \textbf{Twelfth Author\textsuperscript{1}},
%\\
%  \textbf{Thirteenth Author\textsuperscript{3}},
%  \textbf{Fourteenth F. Author\textsuperscript{2,4}},
%  \textbf{Fifteenth Author\textsuperscript{1}},
%  \textbf{Sixteenth Author\textsuperscript{1}},
%\\
%  \textbf{Seventeenth S. Author\textsuperscript{4,5}},
%  \textbf{Eighteenth Author\textsuperscript{3,4}},
%  \textbf{Nineteenth N. Author\textsuperscript{2,5}},
%  \textbf{Twentieth Author\textsuperscript{1}}
%\\
%\\
%  \textsuperscript{1}Affiliation 1,
%  \textsuperscript{2}Affiliation 2,
%  \textsuperscript{3}Affiliation 3,
%  \textsuperscript{4}Affiliation 4,
%  \textsuperscript{5}Affiliation 5
%\\
%  \small{
%    \textbf{Correspondence:} \href{mailto:email@domain}{email@domain}
%  }
%}

\begin{document}
\maketitle
\begin{abstract}
We study the failure modes of large language model (LLM) agents on
long-horizon operating system (OS) tasks using the OS domain of
AgentBench, a 144-task suite in which an agent must translate natural-
language instructions into bash commands under an 8-round interaction
cap. We evaluate two proprietary LLM agents, gpt-5.4-nano and
gemini-2.5-flash, under both the original AgentBench prompt
(\texttt{os-std}) and an optimized variant (\texttt{os-std-opt}) that
augments the system message with eight behavioral guidelines targeting
commonly observed agent failures. The optimized prompt yields task-level
accuracy gains of $+11.1$ points for gpt-5.4-nano and $+6.9$ points for
gemini-2.5-flash, but more than half of all 144 tasks still fail under
every configuration. To explain this gap, we introduce a seven-category
automatic failure taxonomy that classifies each failed trajectory by
inspecting its action sequence. Our central finding is that prompt-only
intervention does not monotonically improve agent behavior; instead, it
\emph{redistributes} failures across categories: under-exploration
failures shrink while round-limit exhaustion and repetition loops grow,
with the dominant failure mode being model-specific.
\end{abstract}

\section{Introduction}

Large language models (LLMs) are increasingly deployed not as single-shot
text generators but as \emph{agents} that interact with external
environments over multiple turns to accomplish complex tasks
\citep{yao2022react, liu2023agentbench, zhou2023webarena}. A particularly
demanding setting is the operating system (OS) domain, in which an agent
must translate a natural-language instruction (e.g., ``recursively set
all files to read-only except those owned by the current user'') into a
sequence of bash commands, observe the resulting system output, and
adapt its strategy until the task is complete. Unlike web or QA agents,
OS agents operate in a regime where commands are syntactically rigid,
errors can irreversibly alter system state, and many tasks have a single
correct numeric answer that is graded by exact match.

The OS domain of AgentBench \citep{liu2023agentbench} provides the
canonical testbed for this setting: 144 evaluation tasks spanning
question-answering and state-modifying operations, executed inside
isolated Docker containers under an 8-round interaction cap. Existing
results paint a sobering picture: even frontier proprietary LLMs solve
fewer than half of these tasks under the default prompt, and aggregate
accuracy alone does not explain \emph{why} they fail. A natural
follow-up question is whether inference-only interventions---prompt
engineering, verbal self-reflection \citep{shinn2023reflexion}, or
iterative self-refinement \citep{madaan2023selfrefine}---can close this
gap, and if so, by mitigating which specific failure modes.

Prior work has sketched failure-mode taxonomies for adjacent settings.
ReAct \citep{yao2022react} reports that 47\% of its HotpotQA failures
take the form of repetitive thought--action loops, motivating
self-reflection-based interventions. PLAN-AND-ACT
\citep{erdogan2025plan} diagnoses plan quality, rather than action
grounding, as the bottleneck for current LLM agents on web tasks.
However, no comparable taxonomy exists for OS tasks specifically, and it
is unclear whether interventions that work on QA or web reasoning
transfer to the rigid syntax and irreversible state of shell command
execution.

In this paper we replicate AgentBench-OS with two contemporary
proprietary LLM agents---gpt-5.4-nano and gemini-2.5-flash---and
systematically characterize the failure modes that survive the strongest
inference-only intervention available to us: an optimized system prompt
augmented with eight behavioral guidelines targeting the specific
pitfalls of the AgentBench-OS environment. We make three contributions:

\begin{enumerate}
\setlength\itemsep{2pt}
\item \textbf{A reproducible benchmark replication.} We evaluate
gpt-5.4-nano and gemini-2.5-flash on all 144 OS tasks under both the
original AgentBench prompt (\texttt{os-std}) and our optimized variant
(\texttt{os-std-opt}), yielding 576 task trajectories with full
\texttt{tool\_calls} traces. The optimized prompt improves task-level
accuracy by $+11.1$ points for gpt-5.4-nano and $+6.9$ points for
gemini-2.5-flash.

\item \textbf{An automatic seven-category failure taxonomy.} We classify
each failed trajectory by inspecting its action sequence in priority
order, distinguishing premature commitment (Premature Finish, Snap Wrong
Answer), inadequate exploration (Round-Limit Exhaustion), inadequate
stopping (Repetitive Bash), and answer-format failures (Late Wrong
Answer). The taxonomy is rule-based and reproducible from the structured
log files alone.

\item \textbf{A failure-redistribution finding.} We show that prompt-only
intervention does not monotonically improve agent behavior; instead, it
\emph{redistributes} failures across categories. For gpt-5.4-nano,
under-exploration failures collapse (Premature Finish: $40 \to 2$, Snap
Wrong Answer: $41 \to 9$) but new failure modes emerge (Round-Limit
Exhaustion: $2 \to 10$, Repetitive Bash: $2 \to 18$). For
gemini-2.5-flash, the dominant baseline failure is repetitive bash
loops, which the optimized prompt mitigates ($28 \to 10$), but at the
cost of more round-limit exhaustions ($10 \to 21$). The dominant
failure mode is model-specific, and the marginal gains from prompt
engineering depend critically on which baseline profile the underlying
model exhibits.
\end{enumerate}

\section{Related Work}
% 0.5-1 page

The evaluation and improvement of LLM-based agents on interactive tasks has
advanced along three closely related lines: benchmark construction,
reasoning-and-acting interaction paradigms, and inference-only improvement
strategies based on prompting and memory. However, existing work does not
fully characterize which inference-only interventions mitigate specific
failure modes in operating system task execution.

\paragraph{Benchmarks for web and digital agents.}
\citet{deng2023mind2web} introduced a large-scale generalist web agent
benchmark with 2,350 tasks across 137 real-world websites. Even GPT-4
achieves only approximately 36\% step success rate under in-context
learning, showing that generalization to unseen digital environments remains
difficult. \citet{zhou2023webarena} addressed the limitation of offline web
evaluation by deploying self-hosted, fully functional websites in
reproducible Docker containers and evaluating agents through programmatic
state inspection rather than action-sequence matching. The best GPT-4-based
agent in WebArena achieves only 14.41\% success, far below the 78.24\%
human baseline, indicating that live digital environments remain challenging
even for frontier models. \citet{liu2023agentbench} extends this benchmark
perspective to a broader set of environments, including the operating
system domain that is the focus of our replication.

\paragraph{Reasoning and acting paradigms.}
\citet{yao2022react} introduced ReAct, a reasoning-and-acting framework
that interleaves free-form ``thought'' tokens with environment actions.
This format gives the agent an implicit working memory without parameter
updates and has become a useful conceptual template for later agent
frameworks. However, ReAct also reveals a characteristic failure mode:
repetitive thought--action loops account for 47\% of its failures on
HotpotQA, motivating stronger inference-time interventions such as
self-reflection, memory augmentation, and explicit planning. The
Thought--Action--Observation trajectory format introduced by ReAct serves
as the shared interaction interface on which AgentBench and most
subsequent agent frameworks build, and is also the format under which our
experiments operate.

\paragraph{Prompt-based and memory-augmented agent improvement.}
A complementary line of work improves LLM agents without updating model
parameters. \citet{shinn2023reflexion} propose Reflexion, which augments
ReAct with a verbal self-reflection step: after a failed trajectory, the
agent writes a natural-language critique and stores it in an episodic
memory buffer that conditions later attempts, yielding substantial gains
over ReAct on AlfWorld decision-making and HotpotQA reasoning, and
reaching 91\% Pass@1 on HumanEval (an 11-point absolute improvement over
GPT-4). \citet{madaan2023selfrefine} develop Self-Refine, in which a
single LLM iteratively generates, critiques, and revises its own output
without supervised training or reinforcement learning, achieving roughly
a 20-point absolute improvement averaged across seven generation tasks.
\citet{erdogan2025plan} take a complementary architectural approach with
PLAN-AND-ACT, decomposing the agent into a \textsc{Planner} that produces
structured plans and an \textsc{Executor} that grounds them into actions.
A central finding is that even an untrained \textsc{Executor} improves by
34 percentage points on WebArena-Lite when paired with a high-quality
\textsc{Planner}, suggesting that plan quality, not action grounding, is
often the binding constraint for current LLM-based agents.

\paragraph{Position of our work.}
Together, these works provide the benchmark, interaction, and
inference-time improvement foundations for our study. Building on the OS
domain of \citet{liu2023agentbench}, our work examines not only whether
prompt- and memory-based interventions improve task success, but also
which categories of failure they actually mitigate in shell command
execution. Our analysis reveals that such prompt-only interventions tend
to redistribute failures across categories rather than monotonically
improve agent behavior, motivating more targeted training-time or
architectural interventions for long-horizon OS tasks.




\section{Method}

We replicate the OS-domain evaluation of \citet{liu2023agentbench} and
extend it with (i) a comparison between the original AgentBench prompt
and an optimized prompt variant, and (ii) an automatic seven-category
failure taxonomy that classifies every failed trajectory by inspecting
its action sequence. Our analysis is purely inference-time: we do not
modify model weights or the AgentBench check-script grader.

\subsection{Task Formulation}

Each AgentBench-OS task is a tuple $(q, e_0, c)$ consisting of a natural-
language instruction $q$, an initial Docker environment state $e_0$, and
a deterministic check script $c$. Given $q$, the agent operates in an
8-round interaction loop with three available tools:
\texttt{bash\_action} (execute a shell command in the container),
\texttt{answer\_action} (submit a final string answer), and
\texttt{finish\_action} (declare completion without returning a value,
used for state-modifying tasks). At each round the environment returns
the truncated stdout/stderr of the executed command. A task is scored
as a pass if (a) for QA tasks, the submitted answer string exactly
matches the reference under the check script, or (b) for operation
tasks, the post-trajectory environment state passes the check script.
Trajectories that exceed the 8-round cap or terminate via
\texttt{finish\_action} without an answer are scored as failures.

\subsection{Prompt Variants}
\label{sec:prompt_variants}

We evaluate two prompt configurations that differ only in the system
message; tools, round cap, decoding settings, and grader are identical.

\paragraph{Baseline (\texttt{os-std}).}
We use the original AgentBench OS system prompt, which establishes the
role (``a person interacting with a Linux operating system''), specifies
the tool-calling protocol, requires the agent to ``call exactly one of
the provided tools'' per turn, and instructs that answers should be
``exact and precise'' (e.g., a single number).

\paragraph{Optimized (\texttt{os-std-opt}).}
We append eight additional behavioral guidelines to the baseline system
message under an ``Additional guidelines'' header. The guidelines target
failure modes observed in pilot runs and are added verbatim:
\begin{enumerate}
\setlength\itemsep{1pt}
\item Confirm the existence and actual path of any file or directory
with \texttt{find} or \texttt{ls} before operating on it.
\item Do not use heredoc syntax (\texttt{<<EOF},
\texttt{<<\textquotesingle EOF\textquotesingle}, etc.); write files
using \texttt{printf} or \texttt{echo} redirection.
\item When the task asks for a script, also execute it and submit the
actual output, not merely confirmation that the script was created.
\item Before calling \texttt{answer\_action}, verify the answer with at
least one cross-check command; do not accept a result of $0$ or empty
without confirming the target exists.
\item Do not use \texttt{set -e}, \texttt{set -u}, or
\texttt{set -o pipefail}, which can cause silent script aborts in this
environment.
\item Do not use shell arithmetic expansion \texttt{echo \$((expr))},
which returns empty output in this environment; use
\texttt{python3 -c} or \texttt{expr} instead.
\item Commit to a consistent result: if two independent methods agree,
call \texttt{answer\_action} immediately, and never re-run the same or
equivalent query more than twice.
\item The submitted answer must be a plain value only, with no
equations, labels, or intermediate steps (e.g., submit
\texttt{100}, not \texttt{4+13+83=100}).
\end{enumerate}

These guidelines are designed to mitigate distinct failure modes:
guidelines~1, 4, and 7 target premature commitment and snap wrong
answers; guideline~7 also constrains repetitive bash loops; guideline~8
addresses format-induced wrong answers (e.g., trajectories that compute
the right number but submit it in an ungraded form); and
guidelines~2, 5, and 6 address shell-execution pitfalls specific to the
AgentBench Docker environment.

\subsection{Failure-Mode Classification}
\label{sec:method_taxonomy}

For each failed trajectory we apply an automatic rule-based classifier
that assigns exactly one of seven categories by inspecting the agent's
action sequence in priority order: \emph{Premature Finish},
\emph{Snap Wrong Answer}, \emph{Late Wrong Answer},
\emph{Round-Limit Exhaustion}, \emph{Repetitive Bash},
\emph{Truncation Handling}, and \emph{Other}. The full category
definitions and the rationale for the priority order are given in
Section~\ref{sec:failure_taxonomy}. The classifier operates on the
structured \texttt{tool\_calls} sequence stored in each
\texttt{runs.jsonl} record, which makes it deterministic and
reproducible from the trajectory logs alone---no human labeling is
required.

\subsection{Models and Inference Settings}

We evaluate two proprietary LLM agents accessed via their respective
chat APIs: \textbf{gpt-5.4-nano} (an OpenAI ChatCompletion endpoint) and
\textbf{gemini-2.5-flash} (a Google GenAI endpoint). Both models support
function-calling, which is required for the AgentBench tool-call
protocol, and are queried with deterministic decoding settings. Each
model is run on all 144 OS tasks under both prompt variants, yielding
four configurations and a total of $4 \times 144 = 576$ task
trajectories.


\section{Experiment}
\label{sec:experiment}

To better understand the characteristics and difficulty of the Operating System (OS) domain in AgentBench, we perform an exploratory analysis of the dataset before conducting experiments. This analysis helps us understand the composition of tasks and the potential challenges faced by LLM-based agents in interactive environments.

\subsection{Dataset Overview}

The OS domain in AgentBench contains 144 evaluation tasks designed to test whether an agent can interact with a Unix-like environment using executable shell commands. Each task provides a natural language instruction that describes a system-related objective, such as retrieving information from the system or modifying the system state. The agent must translate this instruction into a sequence of valid shell commands that can be executed in the environment.

Unlike traditional NLP benchmarks, these tasks require interaction with a real operating system environment. The agent must generate commands, observe execution results, and adapt its behavior based on the feedback returned by the system. This interactive setting makes the benchmark suitable for evaluating LLM-based agents that operate in real-world environments.

\subsection{Task Types}

The dataset includes two main types of tasks: \textit{question-answering (QA)} tasks and \textit{operation} tasks.

QA tasks require the agent to retrieve information from the operating system environment and return the correct answer. For example, the agent may need to search files or inspect system directories to identify specific information.

Operation tasks, on the other hand, require the agent to modify the system state. These tasks involve commands that change files, directories, or permissions. Because incorrect commands may alter the environment state, these tasks require both accurate command generation and careful reasoning about system behavior.

The coexistence of these two task types allows the benchmark to evaluate both information retrieval ability and system manipulation capability of LLM-based agents.




To reproduce the Operating System (OS) evaluation proposed in AgentBench (Liu et al., 2024), we design an experimental protocol that evaluates whether an LLM-based agent can complete operating system tasks through interactive command execution. Our experimental setup follows the original benchmark design, including the execution environment, interaction procedure, prompt construction strategy, and evaluation metrics.

\subsection{Environment Setup}

All experiments are conducted in the Operating System (OS) environment defined in AgentBench. In this setting, the agent interacts with a real bash command-line interface, where generated commands can be executed directly by the system. The environment is implemented using Docker containers, allowing each task to run in an isolated system instance to ensure reproducibility and prevent interference between tasks.

The OS benchmark includes tasks that require the agent to interact with system files, directories, and permissions. Some tasks require retrieving information from the system, while others require modifying the system state through command execution. According to the benchmark design, the agent must generate shell commands that can be executed in the environment to solve the task.

The OS domain contains two types of tasks.
Question Answering (QA) tasks require the agent to generate commands that retrieve specific information from the system and produce a final answer.
Operation tasks require the agent to execute commands that modify the system state, such as updating files, directories, or user permissions. These two task types allow the benchmark to evaluate both information retrieval abilities and system manipulation capabilities of LLM-based agents.

Before a task begins, the environment may initialize specific system states such as directory structures, file permissions, or user settings. This controlled setup allows the benchmark to test whether the agent can correctly generate executable shell commands that solve the given task. The environment then executes the generated commands and returns the resulting system outputs to the agent as feedback.

\subsection{Interaction Protocol}

The evaluation follows a multi-turn interaction process between the agent and the environment. The user first provides a natural language instruction describing the OS task. The agent then produces an action in the form of a shell command.

Once the command is executed, the environment returns the execution result to the agent as an observation. Based on this feedback, the agent decides the next action. This process continues until either the task is successfully completed or the interaction reaches the predefined maximum number of rounds.

AgentBench models this evaluation setting as an interactive agent environment, where the agent repeatedly observes the system state and produces actions until the task terminates.



\section{Evaluation}

We evaluate two proprietary LLM agents, \textit{gpt-5.4-nano} and \textit{gemini-2.5-flash}, on all 144 OS tasks under two prompt variants: the baseline AgentBench prompt (\texttt{os-std}) and an optimized variant (\texttt{os-std-opt}) that augments the system message with eight behavioral guidelines targeting the failure modes observed in pilot runs (e.g.\ ``commit to a consistent result'', ``never re-run the same query more than twice'', ``return a plain numeric answer''). All experiments use a maximum of 8 interaction rounds and the AgentBench check-script grader; details are reported in Section~\ref{sec:experiment}.

\subsection{Overall Accuracy}

Table~\ref{tab:accuracy} reports task-level accuracy. The optimized prompt yields a consistent gain across both models: $+11.1$ points for gpt-5.4-nano and $+6.9$ points for gemini-2.5-flash. However, more than half of the 144 tasks are still failed by every configuration, indicating that prompt-only intervention is not sufficient to close the long-horizon OS gap.

\begin{table}[t]
\centering
\small
\begin{tabular}{lrrr}
\toprule
Model & \texttt{os-std} & \texttt{os-std-opt} & $\Delta$ \\
\midrule
gpt-5.4-nano       & 35.4\% & 46.5\% & $+11.1$ \\
gemini-2.5-flash   & 42.4\% & 49.3\% & $+6.9$  \\
\bottomrule
\end{tabular}
\caption{Task-level accuracy on the 144-task AgentBench-OS suite. \texttt{os-std} is the baseline AgentBench prompt; \texttt{os-std-opt} adds eight behavioral guidelines.}
\label{tab:accuracy}
\end{table}

\begin{figure}[t]
\centering
\begin{tikzpicture}
\begin{groupplot}[
  group style={group size=1 by 2, vertical sep=1.05cm},
  width=\columnwidth,
  height=4.4cm,
  ybar,
  bar width=4.5pt,
  ymin=0,
  enlarge x limits=0.08,
  xtick={1,2,3,4,5,6,7,8},
  ylabel={\# Tasks},
  ylabel style={font=\small},
  tick label style={font=\footnotesize},
  legend style={font=\footnotesize, draw=none, fill=none,
                /tikz/every even column/.append style={column sep=4pt}},
  title style={font=\small, yshift=-3pt},
  every axis plot/.append style={draw=white, line width=0.3pt},
]

\nextgroupplot[
  title={gpt-5.4-nano \quad mean rounds: std=2.79, opt=5.55},
  ymax=120,
  legend pos=north east,
  nodes near coords,
  nodes near coords style={font=\tiny, /pgf/number format/precision=0,
                           /pgf/number format/fixed},
]
\addplot[fill=stdblue]    coordinates {(1,1)(2,101)(3,15)(4,5)(5,11)(6,4)(7,1)(8,6)};
\addplot[fill=optorange]  coordinates {(1,0)(2,16)(3,16)(4,16)(5,22)(6,17)(7,13)(8,44)};
\legend{std, opt}

\nextgroupplot[
  title={gemini-2.5-flash \quad mean rounds: std=4.33, opt=4.57},
  ymax=72,
  xlabel={Number of interaction rounds},
  xlabel style={font=\small},
  legend pos=north east,
  nodes near coords,
  nodes near coords style={font=\tiny, /pgf/number format/precision=0,
                           /pgf/number format/fixed},
]
\addplot[fill=stdblue]    coordinates {(1,1)(2,61)(3,20)(4,7)(5,4)(6,5)(7,5)(8,41)};
\addplot[fill=optorange]  coordinates {(1,0)(2,37)(3,30)(4,13)(5,20)(6,3)(7,4)(8,37)};
\legend{std, opt}

\end{groupplot}
\end{tikzpicture}
\caption{Trajectory length distribution under the baseline (\texttt{os-std})
versus optimized (\texttt{os-std-opt}) prompt for both models. The optimized
prompt does not uniformly extend trajectories: gpt-5.4-nano shifts mass
dramatically out of round~2 ($101 \rightarrow 16$) and into the 8-round cap
($6 \rightarrow 44$), nearly doubling its mean trajectory length
($2.79 \rightarrow 5.55$). gemini-2.5-flash, which already exhibited a
bimodal distribution under the baseline, shows a smaller redistribution
(mean $4.33 \rightarrow 4.57$). The new mass at round~8 corresponds to
round-limit exhaustion failures, which the taxonomy in
Section~\ref{sec:failure_taxonomy} quantifies.}
\label{fig:rounds}
\end{figure}

\subsection{Failure-Mode Taxonomy}
\label{sec:failure_taxonomy}

To go beyond aggregate accuracy, we automatically classify every failed trajectory by inspecting its action sequence. Each failed run is assigned to exactly one of seven mutually exclusive categories:

\begin{itemize}
\setlength\itemsep{2pt}
\item \textbf{Premature finish (PF)}: agent calls \texttt{finish\_action} without ever returning an answer, even though the task has a definite answer.
\item \textbf{Snap wrong answer (SWA)}: agent calls \texttt{answer\_action} within $\le 3$ rounds with an incorrect value, indicating insufficient exploration.
\item \textbf{Late wrong answer (LWA)}: agent calls \texttt{answer\_action} after $>3$ rounds with an incorrect value (typically wrong reasoning or wrong output format despite extensive interaction).
\item \textbf{Round-limit exhaustion (RLE)}: trajectory hits the 8-round cap with no \texttt{answer\_action} or \texttt{finish\_action} ever issued.
\item \textbf{Repetitive bash (RB)}: the agent issues three or more effectively identical \texttt{bash\_action} calls in a row, regardless of how the trajectory ends.
\item \textbf{Truncation handling (TH)}: at least one tool message contains a truncation notice and the agent fails to recover.
\item \textbf{Other}: residual cases not matching any of the above.
\end{itemize}

Table~\ref{tab:failure_modes} reports the per-category counts for both models and both prompt variants.

\begin{table*}[t]
\centering
\small
\begin{tabular}{l@{\hskip 12pt}rrr@{\hskip 18pt}rrr}
\toprule
& \multicolumn{3}{c}{\textbf{gpt-5.4-nano}} & \multicolumn{3}{c}{\textbf{gemini-2.5-flash}} \\
\cmidrule(lr){2-4}\cmidrule(lr){5-7}
Failure category & \texttt{std} & \texttt{opt} & $\Delta$ & \texttt{std} & \texttt{opt} & $\Delta$ \\
\midrule
Premature finish (PF)         & 40 & 2  & $-38$ & 9  & 5  & $-4$  \\
Snap wrong answer (SWA)        & 41 & 9  & $-32$ & 27 & 22 & $-5$  \\
Late wrong answer (LWA)        & 3  & 30 & $+27$ & 7  & 13 & $+6$  \\
Round-limit exhaustion (RLE)   & 2  & 10 & $+8$  & 10 & 21 & $+11$ \\
Repetitive bash (RB)           & 2  & 18 & $+16$ & 28 & 10 & $-18$ \\
Truncation handling (TH)       & 2  & 7  & $+5$  & 1  & 2  & $+1$  \\
Other                          & 3  & 1  & $-2$  & 1  & 0  & $-1$  \\
\midrule
Total failed                   & 93 & 77 & $-16$ & 83 & 73 & $-10$ \\
\bottomrule
\end{tabular}
\caption{Failure-mode breakdown across the four configurations (144 tasks each). Each failed trajectory is assigned to exactly one category in priority order. Negative $\Delta$ indicates fewer failures of that type under the optimized prompt.}
\label{tab:failure_modes}
\end{table*}

\subsection{What the Improved Prompt Mitigates}

As Figure~\ref{fig:rounds} visualizes at the trajectory-length level, the two models exhibit qualitatively different baseline failure profiles, and the optimized prompt acts on each profile differently.

\paragraph{gpt-5.4-nano: under-exploration is the dominant failure.}
Under \texttt{os-std}, $87\%$ of failures (81/93) are accounted for by just two categories, both reflecting under-exploration: the agent either calls \texttt{finish\_action} after a single bash call (PF, $n{=}40$) or commits to an answer in the first 1--2 rounds (SWA, $n{=}41$). The model averaged only $2.79$ rounds per trajectory and $101/144$ tasks terminated by round~2. Adding the eight behavioral guidelines collapses these two categories simultaneously (PF: $40\to 2$, SWA: $41\to 9$), accounting for the entire $+11.1$-point accuracy gain.

\paragraph{gemini-2.5-flash: repetition is the dominant failure.}
Under \texttt{os-std}, gemini's failures are dominated instead by repetitive bash loops (RB, $n{=}28$, $34\%$ of failures) together with snap wrong answers (SWA, $n{=}27$). The optimized prompt's explicit ``never re-run the same or equivalent query more than twice'' rule directly cuts repetition (RB: $28\to 10$). The improvement is smaller in absolute terms ($+6.9$ points) but targets a different failure mode than for gpt.

\subsection{Remaining and Newly Induced Problems}

A central finding of this analysis is that prompt-only intervention does not strictly improve agent behavior: it \emph{redistributes} failures across categories. Three new or amplified failure modes are visible in Table~\ref{tab:failure_modes}.

\paragraph{(1) Late wrong answers replace snap wrong answers.}
For both models, \texttt{LWA} grows substantially under \texttt{os-std-opt} ($+27$ for gpt, $+6$ for gemini). The agent is now willing to explore for longer, but a non-trivial fraction of those longer trajectories still terminate with an incorrect \texttt{answer\_action}. Inspection of these traces shows they are often format failures rather than reasoning failures: e.g.\ task~143 (compute total folder size) is solved with a correct \texttt{du} pipeline but the answer string fails the grader's exact-match rule.

\paragraph{(2) Round-limit exhaustion grows for both models.}
\texttt{RLE} increases from $2{\to}10$ for gpt and from $10{\to}21$ for gemini. Encouraging the agent to commit later trades premature termination for hitting the 8-round cap. A representative case is gpt-opt task~140, where the agent spends 8 rounds writing and rewriting variations of the same \texttt{awk} expression for a log-counting task and never produces \texttt{answer\_action}.

\paragraph{(3) gpt-5.4-nano develops repetitive bash loops it did not previously have.}
While gemini's repetition rate falls under the optimized prompt, gpt-5.4-nano's rises sharply (RB: $2\to 18$). When deprived of its default ``finish quickly'' shortcut, gpt frequently regresses into re-issuing slight variants of the same command, suggesting that the model lacks an internal stopping signal once forced to keep exploring. This pattern echoes the well-documented thought--action loop failure mode on HotpotQA reported for ReAct~\citep{yao2022react}.

\paragraph{Hard-task ceiling.}
Across all four configurations, $33$--$39$ tasks fail at the 8-round cap. Manual inspection of these trajectories suggests that for a substantial subset, the bottleneck is not exploration depth but plan quality: tasks requiring composing a multi-step shell pipeline, generating prime-number filename mappings, or reasoning about file-system semantics beyond simple \texttt{find}/\texttt{wc} idioms are not solved by either model under either prompt. This residual ceiling motivates evaluating planner-decomposed agents in future work.

\paragraph{Generality across model families.}
Our current evaluation uses two proprietary, closed-weights models. To assess whether the failure-redistribution pattern is intrinsic to the OS domain or specific to these systems, we plan to extend the evaluation to open-weights models in the same parameter range. The taxonomy in Table~\ref{tab:failure_modes} is designed to accept additional rows without modification.

\section{Conclusion}

We replicated the AgentBench-OS evaluation with two contemporary
proprietary LLMs and analyzed the impact of an optimized system prompt
augmented with eight behavioral guidelines. The optimized prompt yields
consistent task-level accuracy gains over the AgentBench baseline
($+11.1$ points for gpt-5.4-nano and $+6.9$ points for
gemini-2.5-flash), but more than half of the 144 tasks remain unsolved
under every configuration we tested. Beyond aggregate accuracy, our
seven-category automatic failure taxonomy reveals that prompt-only
interventions do not monotonically improve agent behavior; instead,
they \emph{redistribute} failures across categories, with the specific
redistribution being highly model-dependent. gpt-5.4-nano's dominant
baseline failure is under-exploration, which the optimized prompt
collapses---but at the cost of new round-limit and repetition
failures. gemini-2.5-flash's dominant baseline failure is repetitive
bash loops, which the optimized prompt mitigates with smaller
side-effects. The presence of a hard-task ceiling ($33$--$39$ tasks
fail at the round cap across all four configurations) suggests that no
amount of prompt engineering will close the remaining gap, and that
planner-decomposed agents \citep{erdogan2025plan} or environment-grounded
training are necessary next steps.

\paragraph{Limitations.}
Our study has three principal limitations. First, both evaluated models
are proprietary and closed-weights; whether the failure-redistribution
pattern is intrinsic to the OS domain or specific to these systems
remains an empirical question that requires evaluation on open-weights
models in the same parameter range. Second, our optimized prompt was
hand-tuned in a small pilot study and is not necessarily Pareto-optimal;
automated prompt search may yield further gains, and the specific
guidelines we use may interact with idiosyncratic features of the
AgentBench Docker environment. Third, our 8-round interaction cap
follows the AgentBench convention but is short relative to the genuine
difficulty of some hard-task tail cases; the round-limit-exhaustion
category may partially reflect this cap rather than agent capability.



\bibliography{custom}

\appendix

% \section{Example Appendix}
% \label{sec:appendix}

% This is an appendix.

\end{document}
